\documentclass[12pt]{article}
\usepackage[utf8]{inputenc}
\usepackage[T1]{fontenc}
\usepackage[margin=1in]{geometry}
\usepackage{setspace}
\usepackage{xcolor}
\usepackage[colorlinks=true, urlcolor=blue!60!black, linkcolor=blue!60!black]{hyperref}
\setlength{\parindent}{2em}
\setlength{\parskip}{1.25em}
\renewcommand{\baselinestretch}{1.0}
\begin{document}
\begin{center}
{\LARGE \textit{United States v. Curtiss-Wright Export Corp.}, 299 U.S. 304 (1936)\par}
\vspace{0.8em}
{\large Case Brief --- Four-Page Edition\par}
\end{center}
\vspace{1em}

\section{Memory Jogger}

A company indicted for selling machine guns to Bolivia in violation of a presidentially activated Chaco War arms embargo lost its nondelegation challenge: Congress may give the President far broader discretion in foreign affairs than in domestic matters, in an opinion famous for calling the President the ``sole organ'' of the nation's external relations (\href{https://supreme.justia.com/cases/federal/us/299/304/}{299 U.S. 304}).

\section{Detailed Case Facts}

During the Chaco War between Bolivia and Paraguay, Congress passed a Joint Resolution on May 28, 1934. If the President found that barring arms sales to the belligerents ``may contribute to the reestablishment of peace between those countries,'' and, after consulting other American Republics, so proclaimed, selling arms to them became a crime punishable by up to \$10,000 and two years in prison, ``except under such limitations and exceptions as the President prescribes.'' President Roosevelt issued the proclamation that day and revoked it in November 1935, preserving liabilities already incurred.

\section{Procedural History}

In January 1936, a grand jury in the Southern District of New York indicted Curtiss-Wright Export Corp. for conspiring to sell 15 machine guns to Bolivia. The defendants demurred, arguing that the Resolution unconstitutionally delegated legislative power (citing \textit{Panama Refining} and \textit{Schechter Poultry}), that the proclamation lacked required findings, and that revocation extinguished liability. The district court sustained the demurrer on the delegation ground only. The United States appealed directly under the Criminal Appeals Act.

\section{Judicial Votes}

Reversed and remanded, \textbf{7--1} (decided December 21, 1936).
\begin{itemize}
\item \textbf{Majority:} Justice Sutherland, joined by Chief Justice Hughes and Justices Van Devanter, Brandeis, Butler, Roberts, and Cardozo.
\item \textbf{Concurrences:} None.
\item \textbf{Dissent:} Justice McReynolds, without opinion. Justice Stone did not participate.
\end{itemize}

\section{Holding}

The Joint Resolution was not an unconstitutional delegation, because Congress may confer far broader discretion on the President in external affairs than in purely domestic matters. The proclamation tracked the statute and satisfied every jurisdictional requirement, and revocation did not repeal the Resolution or extinguish liability for offenses committed while the embargo was in force.

\section{Analysis of Opinions}

\textbf{Majority (Sutherland).} Assuming the delegation might fail in domestic affairs, Sutherland stressed that the Resolution's ``whole aim'' was ``to affect a situation entirely external to the United States.'' First, enumerated powers are true ``only in respect of our internal affairs''; because the States ``severally never possessed international powers,'' external sovereignty passed from the Crown to the United States collectively and did not depend on the Constitution's affirmative grants. Second, quoting John Marshall's 1800 House speech, Sutherland called the President ``the sole organ of the nation in its external relations,'' with confidential sources, secrecy, and dispatch that make narrow domestic-style standards unwise. Third, a ``uniform, long-continued and undisputed legislative practice'' of embargo statutes since 1794 confirmed the point. The proclamation's recitals sufficed, and revocation ended the embargo only prospectively.

\textbf{McReynolds (dissenting).} He would have affirmed the district court.

\textbf{Later cases.} Justice Jackson's \textit{Youngstown} concurrence treated the case as one of action under an Act of Congress and much of its language as dicta. \textit{Zivotofsky v. Kerry} (2015) repudiated the broad ``sole organ'' dicta. In \textit{Learning Resources v. Trump} (2026), the Court invalidated IEEPA tariffs; Justice Gorsuch cabined \textit{Curtiss-Wright} to a targeted embargo, while Justices Kavanaugh and Thomas relied on it in dissent.

\section{Examples: Future Applications}

\begin{enumerate}
\item \textbf{Same-Side (Upheld): Chip Export Ban.} A statute letting the President ban advanced-chip exports upon a finding of a strategic threat is strong \textit{Curtiss-Wright} territory.
\item \textbf{Same-Side (Upheld): Sanctions on an Aggressor.} A ban on aircraft-part sales to a government that invades its neighbor, with presidential exceptions, nearly copies the 1934 Resolution.
\item \textbf{Opposite-Side (Invalid): Export Fee.} A 20\% fee on chip exports runs into the Export Clause; foreign-affairs power ``must be exercised in subordination to the applicable provisions of the Constitution.''
\item \textbf{Opposite-Side (Invalid): Order to State Courts.} A presidential memorandum, with no statute, directing state courts to reopen convictions fails under \textit{Youngstown} and \textit{Medell\'in v. Texas}.
\item \textbf{Fence-Sitter: Foreign-Owned App.} Banning a foreign-controlled app rests on confidential intelligence but operates on domestic commerce and speech.
\end{enumerate}

\section{Critique}

\textbf{Flawed history.} Charles Lofgren and David Levitan showed that States exercised external powers between 1776 and 1789 and that the Constitution divides foreign-affairs powers between Congress and the President.

\textbf{Misused Marshall and a logical leap.} Marshall described the President as the sole \textit{communicator} of policy already set by law or treaty, not its sole \textit{maker}. And even if external sovereignty passed to the United States as a whole, that explains federal rather than state power, not why it belongs to the President rather than Congress.

\section{Key Quotations}

\begin{quote}
``It is important to bear in mind that we are here dealing not alone with an authority vested in the President by an exertion of legislative power, but with such an authority plus the very delicate, plenary and exclusive power of the President as the sole organ of the federal government in the field of international relations---a power which does not require as a basis for its exercise an act of Congress but which, of course, like every other governmental power, must be exercised in subordination to the applicable provisions of the Constitution.''
\end{quote}
\noindent --- Justice Sutherland, 299 U.S. at 319--320.

\begin{quote}
``[C]ongressional legislation which is to be made effective through negotiation and inquiry within the international field must often accord to the President a degree of discretion and freedom from statutory restriction which would not be admissible were domestic affairs alone involved.''
\end{quote}
\noindent --- Justice Sutherland, 299 U.S. at 320.

\begin{quote}
``In this vast external realm, with its important, complicated, delicate and manifold problems, the President alone has the power to speak or listen as a representative of the nation. He makes treaties with the advice and consent of the Senate; but he alone negotiates.''
\end{quote}
\noindent --- Justice Sutherland, 299 U.S. at 319.

\end{document}
